\documentclass[11pt,a4paper]{article}
\usepackage[utf8]{inputenc}
\usepackage[margin=1in]{geometry}
\usepackage{amsmath,amssymb}
\usepackage{booktabs}
\usepackage{xcolor}
\usepackage{hyperref}
\usepackage{fancyhdr}
\usepackage{enumitem}

\setlength{\headheight}{14pt}

\hypersetup{
    colorlinks=true,
    linkcolor=blue,
    filecolor=magenta,      
    urlcolor=cyan,
    pdftitle={Project Report: Simple Joule Thief Circuit},
    pdfpagemode=UseOutlines,
}

\pagestyle{fancy}
\fancyhf{}
\rhead{Joule Thief Voltage Booster}
\lhead{Technical Project Report}
\rfoot{Page \thepage}

\title{\textbf{\Huge Engineering Project Report}\\[0.5em] \Large Simple Joule Thief Circuit: Low-Power Inductive Voltage Booster}
\author{\textbf{Bigyanlabs Engineering Team}\\ Energy Harvesting \& Power Electronics Division}
\date{\today}

\begin{document}

\maketitle
\thispagestyle{empty}

\begin{abstract}
This report documents the design, electromagnetic analysis, and laboratory evaluation of a Joule Thief self-oscillating voltage booster circuit. The Joule Thief is a minimalist flyback boost converter designed to extract residual energy from depleted chemical batteries ($0.8\,\text{V} - 1.1\,\text{V}$). By utilizing a dual-winding (bifilar) toroidal inductor, an NPN switching transistor (2N2222), and a single base resistor, the circuit generates high-frequency voltage spikes ($3\,\text{V} - 15\,\text{V}$ at $50\,\text{kHz} - 500\,\text{kHz}$) capable of driving standard LEDs ($V_F \approx 3.0\,\text{V}$). This document presents the five-phase self-oscillation cycle, inductor saturation mechanics, winding rules, component selection criteria, and practical performance limitations.
\end{abstract}

\newpage
\tableofcontents
\newpage

\section{Introduction \& Operating Concept}
Conventional light-emitting diodes (LEDs) require a minimum forward threshold voltage ($V_F \approx 2.0\,\text{V} - 3.4\,\text{V}$) to initiate photon emission. Standard $1.5\,\text{V}$ AA/AAA alkaline cells are typically discarded once their output voltage drops below $1.1\,\text{V}$, even though significant chemical energy remains unextracted.

The **Joule Thief** overcomes this voltage threshold by using inductive energy storage and rapid transistor switching, boosting input voltages as low as $0.7\,\text{V} - 0.8\,\text{V}$ into high-voltage pulses.

\section{Bill of Materials}
\begin{table}[h!]
\centering
\caption{Bill of Materials for Joule Thief Circuit}
\vspace{0.5em}
\begin{tabular}{@{}c l c l p{6.5cm}@{}}
\toprule
\textbf{S.No.} & \textbf{Component} & \textbf{Qty} & \textbf{Part Number} & \textbf{Circuit Function} \\
\midrule
1 & NPN Transistor & 1 & 2N2222 / 2N3904 & High-speed switching device \\
2 & Toroidal Core & 1 & Ferrite Ring Core & Magnetic flux storage medium \\
3 & Magnet Wire & 2 Strands & 28 AWG Enamelled & Primary \& Feedback bifilar windings ($10-15$ turns) \\
4 & Base Resistor & 1 & $1\,\text{k}\Omega$, 1/4W & Limits base drive current and sets frequency \\
5 & Output Load & 1 & 5mm White LED & Light output indicator ($V_F \approx 3.2\,\text{V}$) \\
6 & Power Source & 1 & $1.5\,\text{V}$ AA Battery & Depleted cell ($0.8\,\text{V} - 1.2\,\text{V}$) \\
\bottomrule
\end{tabular}
\end{table}

\section{Toroid Winding Architecture}
Proper phase alignment of the two windings is critical for positive feedback:
\begin{enumerate}
    \item Wind two insulated copper wires (Red = Primary, Yellow = Feedback) simultaneously around the ferrite toroid for $10-15$ turns (**Bifilar Winding**).
    \item Connect the finish of the Primary winding ($R_2$) to the start of the Feedback winding ($Y_1$). This junction forms the **Center Tap** connected to $+V_{\text{battery}}$.
    \item Connect the start of the Primary winding ($R_1$) to the transistor **Collector**.
    \item Connect the finish of the Feedback winding ($Y_2$) through the $1\,\text{k}\Omega$ resistor to the transistor **Base**.
\end{enumerate}

\section{Five-Phase Oscillation Dynamics}

\subsection{Phase 1: Initial Base Conduction}
Current flows from the battery through the Feedback winding and $1\,\text{k}\Omega$ resistor into the Base, turning the transistor slightly ON. Current begins flowing through the Primary winding into the Collector.

\subsection{Phase 2: Regenerative Positive Feedback}
Primary current induces a positive voltage in the Feedback winding via mutual inductance. This increases Base current, rapidly saturating the transistor into full conduction.

\subsection{Phase 3: Magnetic Core Saturation}
Primary current increases until the ferrite core reaches magnetic saturation ($\frac{d\Phi}{dt} \rightarrow 0$). Induced voltage in the feedback winding collapses to zero, cutting off base current and turning the transistor OFF abruptly.

\subsection{Phase 4: Inductive Flyback Voltage Spike}
With the transistor OFF, the collapsing magnetic field generates a high-voltage back-EMF pulse according to Faraday's Law:
\begin{equation}
V_L = -L \frac{di}{dt}
\end{equation}
This voltage spike ($3\,\text{V} - 15\,\text{V}$) adds to the battery voltage, exceeding the LED forward voltage threshold ($V_F$) and lighting the LED.

\subsection{Phase 5: Cycle Reset}
Once magnetic energy is depleted into the LED, the coil voltage drops to zero. Small battery current flows back into the base, initiating Phase 1 again at $50\,\text{kHz} - 500\,\text{kHz}$.

\section{Performance Summary \& Applications}
\begin{itemize}
    \item \textbf{Minimum Input Voltage:} Runs reliably down to $\sim 0.75\,\text{V}$.
    \item \textbf{Oscillation Frequency:} $50\,\text{kHz} - 500\,\text{kHz}$ depending on core permeability and turns ratio.
    \item \textbf{Applications:} Emergency flashlights from dead batteries, micro-energy harvesting, and inductive switching educational kits.
\end{itemize}

\section{Conclusion}
The Joule Thief provides a simple, self-oscillating demonstration of energy extraction from low-voltage sources, transforming weak battery energy into bright LED illumination via flyback induction.

\end{document}
